% Generated draft: requires mathematical and source review.
\documentclass[11pt]{article}
\usepackage[margin=25mm]{geometry}
\usepackage{fontspec}
\setmainfont{Latin Modern Roman}
\usepackage{amsmath,amssymb,mathtools,unicode-math}
\setmathfont{Latin Modern Math}
\usepackage{enumitem,needspace,xurl,hyperref}
\setlength{\parindent}{0pt}
\setlength{\parskip}{5pt}
\setlength{\emergencystretch}{3em}
\newcommand{\N}{\mathbb{N}}
\newcommand{\Z}{\mathbb{Z}}
\newcommand{\Q}{\mathbb{Q}}
\newcommand{\R}{\mathbb{R}}
\newcommand{\C}{\mathbb{C}}
\newcommand{\F}{\mathbb{F}}
\newcommand{\A}{\mathbb{A}}
\newcommand{\PP}{\mathbb{P}}
\newcommand{\E}{\mathbb{E}}
\newcommand{\eps}{\varepsilon}
\newcommand{\dd}{\,\mathrm{d}}
\begin{document}
\section*{UnsolvedMath ID: 700222}
\textbf{Source problem code:} catalogue entry
\subsection{Definitions and mathematical statement}
A variety of groups \(\mathcal V\) is a class of groups defined by group identities. Let \(F=F_{\mathcal V}(x,y)\) be its relatively free group on two designated generators: for every \(G\in\mathcal V\) and every \(a,b\in G\), there is a unique homomorphism \(F\to G\) sending \(x\) to \(a\) and \(y\) to \(b\).

Here, finite presentability means finite presentability as an abstract group: existence of a presentation with finitely many generators and finitely many defining relators. Interpret the source's term ``multiplicator'' as the Schur multiplier
\[
M(F)=H_2(F,\Z),
\]
where the coefficients have trivial \(F\)-action. Equivalently, if \(L\) is the ordinary free group on \(x,y\) and \(R\) is the kernel of its natural epimorphism onto \(F\), then
\[
M(F)\cong \frac{R\cap[L,L]}{[L,R]}.
\]
Here \([L,L]\) is the commutator subgroup, and \([L,R]\) is the normal subgroup of \(L\) generated by the commutators \([\ell,r]=\ell^{-1}r^{-1}\ell r\), with \(\ell\in L\) and \(r\in R\).

\textbf{Question (G. Baumslag).} Is it true that for every variety of groups \(\mathcal V\),
\[
F_{\mathcal V}(x,y)\text{ is not finitely presented}
\quad\Longrightarrow\quad
M\bigl(F_{\mathcal V}(x,y)\bigr)\text{ is not finitely generated as an abelian group}?
\]
The source does not explicitly define ``multiplicator'' or specify the sense of finite presentation; this formulation uses the conventional interpretations stated above.
\subsection{Short English statement}
Must the relatively free group on two generators in any variety of groups have an infinitely generated Schur multiplier whenever it has no finite presentation as an abstract group? This uses the conventional interpretation of the source's term ``multiplicator.''
\subsection{Sources}
\begin{itemize}
\item Supplied local UnsolvedMath registry record, permanent ID \(700222\), exact source problem code \texttt{catalogue entry}, title ``Kourovka Notebook Problem 4.11.''
\item The supplied record attributes the problem to G. Baumslag, proposed in 1973, and cites \emph{The Kourovka Notebook (Unsolved Problems in Group Theory)}, 4th issue (1973), Problem 4.11, as reproduced in the 21st issue edition (2026), \texttt{arXiv:1401.0300v48}. The supplied collection description names E. I. Khukhro and V. D. Mazurov as editors.
\item The registry supplies the status label ``open'' and an update timestamp of \texttt{2026-10-16T14:00:00Z}. These are supplied registry metadata, not an independently verified assessment of the literature or mathematical status. The cited source has not been independently consulted for this draft.
\end{itemize}
\end{document}
